% =====================================================================
% III-B: cấp độ 2. Ablation ở cuối (Bảng tab:l2-ablation); kết quả chính ở Mục IV-C.
% Nguồn số của tab:l2-ablation: experiments/level2/{val,test}/summary.json (dòng kg_triple, shuffled_edges,
% random_emb, random_kgc) và experiments/level2/val/alpha_beta_sweep{,_kgc}.json (β > 0 tốt nhất qua 3 seed).
\subsection{Cấp độ 2: đồ thị tri thức MKG-W và NativE}
\label{sec:method-l2}

\subsubsection{NativE}
NativE~\cite{zhang2024native} là một mô hình hoàn thiện đồ thị tri thức đa phương thức (MMKGC: dự đoán cạnh còn thiếu bằng
cả cấu trúc lẫn ảnh và chữ của thực thể). Lõi cấu trúc là AdvRelRotatE, một biến thể của RotatE~\cite{sun2019rotate}: mỗi
quan hệ là một phép quay trong không gian phức, và mẫu âm khó được sinh bằng huấn luyện đối kháng. Đặc trưng ảnh và chữ được
trộn vào embedding thực thể với trọng số phụ thuộc quan hệ, nên mô hình chịu được thực thể thiếu ảnh hoặc thiếu chữ. Nhóm
dùng NativE theo hai bước.
\begin{itemize}
\item \textbf{Bước a: tái lập.} Huấn luyện trên MKG-W gốc với đúng mã nguồn (commit \texttt{c72291c}) và siêu tham số của
tác giả (1\,000 epoch, dim 250, margin 4, 128 mẫu âm, lr $10^{-4}$, batch 1\,024). So MRR và Hits@1 với bài báo
(Mục~\ref{sec:l2}).
\item \textbf{Bước b: đưa vào truy vấn.} Huấn luyện trên bản \texttt{MKG-W-rmask} (Mục~\ref{sec:leakage}) với 3 seed, rồi dùng
mô hình đã học làm tín hiệu đồ thị.
\end{itemize}

\subsubsection{Liên kết truy vấn và điểm đồ thị}
Câu truy vấn đã che tên đáp án nhưng vẫn hay nhắc tên thực thể khác (``a rock guitarist who was a member of \emph{Pink
Floyd}''). \texttt{EntityLinker} tìm trong câu mọi nhãn và alias của thực thể trong đồ thị, ưu tiên cụm khớp dài nhất (tối
đa 5 token). Gọi $L$ là tập thực thể khớp được, $c$ là thực thể ứng viên. Có ba kênh điểm đồ thị.
\begin{itemize}
\item \texttt{triple}: số cạnh KGC-train trực tiếp giữa $c$ và các thực thể trong $L$. Kênh này không cần huấn luyện.
\item \texttt{emb}: cosine lớn nhất giữa embedding cấu trúc NativE của $c$ và của các thực thể trong $L$.
\item \texttt{kgc}: điểm bộ ba của chính NativE, lấy giá trị lớn nhất qua mọi $l \in L\setminus\{c\}$, mọi quan hệ (169)
và cả hai chiều $(c, r, l)$, $(l, r, c)$. Khác \texttt{emb} ở chỗ có dùng quan hệ. Code tính lại đúng hàm
\texttt{forward()} của NativE và có test đối chiếu với checkpoint thật.
\end{itemize}
Câu không khớp thực thể nào nhận điểm đồ thị 0, nên thứ hạng của nó giữ nguyên như CLIP. Điểm cuối là
\begin{equation}
\alpha\,z(\text{CLIP}) + (1-\alpha)\,z\big(\beta\,z(N) + (1-\beta)\,z(\texttt{triple})\big),
\end{equation}
với $N$ là \texttt{emb} hoặc \texttt{kgc}. Lưới tìm $\alpha \in \{0; 0{,}05; \dots; 1\}$,
$\beta \in \{0; 0{,}25; 0{,}5; 0{,}75; 1\}$, chọn trên val riêng cho từng seed theo R@1 (hoà thì xét MRR). Với
$\beta = 1$, điểm đồ thị chỉ còn tín hiệu của NativE; với $\beta = 0$, chỉ còn số cạnh.

\textbf{Về GAT trên MKG-W.} MKG-W không có scene graph cho từng ảnh, nên bộ xếp lại GATv2 của cấp~1 không áp dụng được. Ở
bước b, tín hiệu từ đồ thị tri thức (đếm cạnh và NativE) thay vai trò của GNN.

\subsubsection{Lời giải thích}\label{sec:l2-expl}
Lời giải thích của một ứng viên là danh sách thực thể khớp được trong câu và các cạnh KGC-train giữa chúng với ứng viên. Đây
chính là các cạnh đã được đếm vào kênh \texttt{triple}, nên khi $\beta = 0$ lời giải thích chứa đúng toàn bộ điểm đồ thị.

\subsubsection{Ablation}\label{sec:l2-ablation}

\begin{table}[t]
\centering
\caption{Ablation của kênh đồ thị trên MKG-W (892 truy vấn mỗi split). Các dòng dùng 3 seed của NativE; $\alpha,\beta$
lấy từ val.}
\label{tab:l2-ablation}
\footnotesize
\setlength{\tabcolsep}{2.5pt}
\resizebox{\columnwidth}{!}{%
\begin{tabular}{lcccc}
\toprule
 & \multicolumn{2}{c}{Val} & \multicolumn{2}{c}{Test} \\
\cmidrule(lr){2-3}\cmidrule(l){4-5}
Cấu hình & R@1 & MRR & R@1 & MRR \\
\midrule
Dòng chính (đếm cạnh, $\beta=0$) & 32,40 & 41,18 & 29,82 & 39,89 \\
Cạnh xáo trộn (giữ bậc của đầu) & \pmstd{21,82}{0,39} & \pmstd{32,03}{0,23} & \pmstd{21,45}{0,06} & \pmstd{32,02}{0,16} \\
NativE thay bằng nhiễu Gauss & 32,40 & 41,18 & 29,82 & 39,89 \\
Thêm cosine NativE, $\beta>0$ tốt nhất\textsuperscript{*} & 31,61 & --- & --- & --- \\
Thêm điểm NativE, $\beta>0$ tốt nhất\textsuperscript{*} & 30,38 & --- & --- & --- \\
\bottomrule
\end{tabular}}
\par\smallskip
{\scriptsize\raggedright \textsuperscript{*}R@1 val lớn nhất qua mọi $\alpha$, mọi $\beta>0$ và 3 seed; lưới chỉ lưu R@1, và
test không chạy vì val chọn $\beta=0$.\par}
\tabsrc{\file{experiments/level2/val/summary.json}, \file{experiments/level2/test/summary.json},
\file{experiments/level2/val/alpha_beta_sweep.json}, \file{experiments/level2/val/alpha_beta_sweep_kgc.json}}
\end{table}

Bảng~\ref{tab:l2-ablation} kiểm tra hai điều.
\begin{itemize}
\item \textbf{Cạnh phải nối đúng thực thể.} Dòng ``cạnh xáo trộn'' đổi chỗ đuôi của các triple trong cùng một quan hệ, giữ
nguyên bậc của đầu, rồi dùng lại $\alpha,\beta$ của dòng chính. R@1 khi đó thấp hơn CLIP thuần (val 21,82 so với 24,33; test
21,45 so với 24,89). Vì phép xáo giữ nguyên bậc của đầu cạnh, mức tăng của dòng chính đến từ việc cạnh nối đúng thực thể,
không đến từ bậc của thực thể.
\item \textbf{NativE không đóng góp khi đã có kênh đếm cạnh.} Ở cả 3 seed, val chọn $\beta = 0$ cho cả kênh cosine lẫn kênh
\texttt{kgc}; mọi $\beta > 0$ đều cho R@1 val thấp hơn (tốt nhất 31,61 với cosine và 30,38 với \texttt{kgc}, so với 32,40).
Do kênh NativE có trọng số 0, thay embedding của nó bằng nhiễu Gauss cho đúng con số của dòng chính. Kết quả của NativE
dùng riêng ($\beta = 1$) ở Bảng~\ref{tab:l2-results}.
\end{itemize}
